% LandmarkDiff: Anatomically-Conditioned Latent Diffusion for
% Photorealistic Facial Surgery Outcome Prediction
% MICCAI 2026 Submission — Springer LNCS Format

\documentclass[runningheads]{llncs}

\usepackage{graphicx}
\usepackage{amsmath,amssymb}
\usepackage{booktabs}
\usepackage{xcolor}
\usepackage[hidelinks]{hyperref}

% Custom commands
\newcommand{\method}{LandmarkDiff}
\newcommand{\ie}{\textit{i.e.}}
\newcommand{\eg}{\textit{e.g.}}
\newcommand{\etal}{\textit{et al.}}
\newcommand{\wrt}{w.r.t.}
\newcommand{\RR}{\mathbb{R}}

\begin{document}

\title{\method{}: Anatomically-Conditioned Latent Diffusion for\\Photorealistic Facial Surgery Outcome Prediction}

\author{Anonymous MICCAI Submission}
\institute{Anonymous Institution}

\maketitle

%======================================================================
\begin{abstract}
Predicting what a patient will look like after facial surgery from a
single 2D photograph would transform surgical consultation, yet current
approaches either demand expensive 3D imaging (CT/CBCT) or resort to
morphing tools that produce obvious artifacts.  We present \method{}, a
framework that generates photorealistic post-operative predictions by
conditioning latent diffusion on anatomically deformed facial landmark
meshes.  Given a clinical photograph, the pipeline extracts 478 landmarks
via MediaPipe, displaces them with procedure-specific Gaussian RBF fields
calibrated from anthropometric data, and feeds a three-channel conditioning
signal---tessellation mesh, Canny edges, and a feathered surgical mask---into
a ControlNet-augmented Stable Diffusion backbone.  LAB-space skin tone
matching and mask compositing preserve identity outside the surgical region.
We cover four procedures (rhinoplasty, blepharoplasty, rhytidectomy,
orthognathic surgery) and evaluate with FID, LPIPS, NME, SSIM, and ArcFace
identity similarity, stratified by Fitzpatrick skin type.
\end{abstract}

%======================================================================
\section{Introduction}
\label{sec:intro}

When a patient walks into a consultation for rhinoplasty or a facelift,
the surgeon's ability to show them what to expect matters enormously.
Yet the tools available today are surprisingly crude: verbal descriptions,
generic before/after galleries from other patients, or commercial morphing
software (Crisalix~\cite{crisalix}, AEDIT) that stretches pixels with
affine warps.  These morphing tools invariably produce visible distortion
at tissue boundaries and cannot synthesize realistic skin texture where
the surgery would alter anatomy.

The research community has made progress on computational surgical
prediction, but almost exclusively through 3D pipelines requiring
cone-beam CT or structured-light scans~\cite{ma2021fcnet,fang2022acmtnet,lee2025camos}.  A CBCT scanner costs \$100K--\$250K,
exposes patients to radiation, and sits in a hospital radiology suite—not
the outpatient cosmetic clinic where the vast majority of the world's 15+
million annual facial procedures~\cite{isaps2023} are actually performed.
These patients have photographs, not CT volumes.

What is missing, then, is a method that can take a single clinical
photograph—a 2D image from a phone camera—and produce a photorealistic
prediction of the surgical outcome while respecting the specific
biomechanics of the planned procedure.  This means \emph{not} arbitrary
warps, but anatomically grounded deformations (\eg, bilateral alar
narrowing for rhinoplasty), and \emph{not} pixel interpolation, but
genuine synthesis of skin texture, shading, and pore-level detail in the
modified region.  Patient satisfaction correlates strongly with
pre-operative expectation management~\cite{honigman2004review}, making
this a problem worth solving well.

We introduce \method{}, which bridges this gap by pairing dense facial
landmark manipulation with ControlNet-conditioned latent
diffusion~\cite{zhang2023controlnet,rombach2022ldm}.  Starting from a
single photograph, the system extracts 478 facial landmarks via
MediaPipe~\cite{lugaresi2019mediapipe}, deforms them according to
procedure-specific Gaussian RBF fields grounded in anthropometric
data~\cite{singh2010plastic}, renders a three-channel conditioning signal
(tessellation mesh, Canny edges, feathered surgical mask), generates the
prediction through ControlNet/Stable Diffusion~1.5, and composites the
result back into the original photograph with LAB-space skin tone
matching.  The entire pipeline runs on a single GPU.

Our main contributions:
\begin{itemize}
    \item A latent diffusion framework for 2D facial surgery prediction
    conditioned on anatomically deformed landmark meshes, removing
    the need for 3D imaging.
    \item Procedure-specific Gaussian RBF deformation fields for
    rhinoplasty, blepharoplasty, rhytidectomy, and orthognathic surgery,
    calibrated against published anthropometric measurements.
    \item A three-channel conditioning scheme that reuses pre-trained
    CrucibleAI ControlNet weights with minimal additional training.
    \item An evaluation protocol stratified by Fitzpatrick skin type,
    providing transparency on performance across skin tones.
\end{itemize}

%======================================================================
\section{Related Work}
\label{sec:related}

\paragraph{3D-Based Surgical Prediction.}
Most computational work on facial surgery simulation has operated in the
3D domain.  Ma~\etal~\cite{ma2021fcnet} learn a fully convolutional mapping
from pre-operative CBCT bone and soft tissue to post-orthognathic
displacement fields.  ACMT-Net~\cite{fang2022acmtnet} adds cross-modal
attention between CT volumes and surface meshes to simulate facial
appearance changes driven by surgically planned bony movements.
More recently, Lee~\etal~\cite{lee2025camos} use conditional multi-scale
autoregressive modeling for orthognathic facial appearance prediction,
achieving strong geometric accuracy on CBCT data.  These methods
produce geometrically faithful results, but they require infrastructure
that most cosmetic clinics simply do not have: CBCT scanners run
\$100K--\$250K, expose patients to ionizing radiation, and occupy
dedicated imaging suites.

\paragraph{2D Surgery Simulation and Face Recognition.}
Work on surgical face transformation in 2D has historically centered on
recognition rather than synthesis.  Singh~\etal~\cite{singh2010plastic}
assembled the IIITD Plastic Surgery Face Database (900 subjects across 6
procedures) and quantified how surgery degrades face recognition accuracy,
providing the anthropometric displacement statistics we draw on for our
deformation model.  Rathgeb~\etal~\cite{rathgeb2020obstacle} confirmed
with deep networks that this degradation persists even for modern
recognition systems.  On the generative side, efforts have been limited:
PlasticGAN~\cite{chandaliya2022plasticgan} trained a holistic GAN for
post-surgery face synthesis but produced noticeable boundary artifacts.
Tanikawa and Yamashiro~\cite{tanikawa2021soft} used deep learning for
soft tissue prediction in mandibular advancement, and
Jung~\etal~\cite{jung2024rhinoplasty} explored AI-powered rhinoplasty
outcome simulation---but both remain limited to narrow procedure scopes
and lack the photorealism of modern generative models.

\paragraph{Diffusion Models and Spatial Conditioning.}
Latent diffusion models~\cite{rombach2022ldm} moved image synthesis into
compressed latent space, yielding photorealistic outputs at tractable
compute cost.  ControlNet~\cite{zhang2023controlnet} introduced a way to
inject spatial structure---edge maps, skeletons, depth---into a frozen
diffusion backbone through zero-initialized parallel encoder blocks,
preserving the learned generative prior while gaining geometric control.
Of particular relevance is the CrucibleAI variant trained on
MediaPipe face mesh tessellations~\cite{crucibleai_controlnet}: it
conditions directly on the 2,556-edge triangulated wireframe, generating
faces that conform to specified landmark geometry.  We repurpose this
model as our generative backbone, sidestepping the data and compute
burden of training a face-conditioned ControlNet from scratch.

\paragraph{Clinical Domain Gap.}
Models trained on curated face datasets (FFHQ, CelebA-HQ) degrade on
real clinical photographs, which suffer from point-source lighting, color
temperature shifts, fluorescent casts, JPEG artifacts, and lens
distortion.  Daneshjou~\etal~\cite{daneshjou2022disparities} documented
this degradation in dermatology, particularly for darker skin tones.  We
tackle this head-on with a domain randomization
approach~\cite{tobin2017domain}: a pool of eight clinical degradation
augmentations applied stochastically during training to close the
distribution gap before it becomes a problem.

\paragraph{Positioning.}
No prior work, to our knowledge, weds diffusion-based photorealistic
synthesis with anatomically grounded, procedure-specific deformations in
a 2D-only surgical prediction system.  Table~\ref{tab:comparison} makes
the comparison explicit.

\begin{table}[t]
\centering
\caption{Comparison with existing surgical prediction approaches.}
\label{tab:comparison}
\small
\begin{tabular}{@{}lcccc@{}}
\toprule
Method & Input & Generative & Procedure- & Photo- \\
       &       & Model      & Specific   & realistic \\
\midrule
Ma~\etal~\cite{ma2021fcnet}       & 3D CBCT & FCN & \checkmark & -- \\
ACMT-Net~\cite{fang2022acmtnet} & 3D CT+Mesh & Attention & \checkmark & -- \\
Lee~\etal~\cite{lee2025camos}     & 3D CBCT & Autoregressive & \checkmark & -- \\
PlasticGAN~\cite{chandaliya2022plasticgan} & 2D Photo & GAN & -- & Partial \\
Jung~\etal~\cite{jung2024rhinoplasty}     & 2D Photo & AI sim. & Rhino only & Partial \\
Crisalix~\cite{crisalix}        & 2D/3D & Morphing & \checkmark & -- \\
\midrule
\textbf{\method{} (Ours)}       & \textbf{2D Photo} & \textbf{LDM+ControlNet} & \textbf{\checkmark} & \textbf{\checkmark} \\
\bottomrule
\end{tabular}
\end{table}

%======================================================================
\section{Methods}
\label{sec:methods}

\method{} operates in five stages (Fig.~\ref{fig:architecture}): (1)~dense
facial landmark extraction, (2)~procedure-specific anatomical deformation,
(3)~multi-channel conditioning signal generation, (4)~ControlNet-guided
latent diffusion, and (5)~identity-preserving mask compositing.

\subsection{Dense Facial Landmark Extraction}
\label{sec:landmarks}

We extract $N = 478$ facial landmarks using MediaPipe Face
Mesh~\cite{lugaresi2019mediapipe}, including 468 face surface points and
10 iris landmarks (5 per eye).  For each input image $I \in \RR^{H
\times W \times 3}$, the detector produces normalized coordinates
$\mathbf{L} = \{(x_i, y_i, z_i)\}_{i=1}^{N}$ with $x_i, y_i \in [0,1]$
representing fractional image position and $z_i$ encoding relative depth.
Landmarks are organized into anatomically meaningful regions—jawline (36
points), orbital contours (16 per eye), nasal dorsum and alar base (25
points), lip vermilion (22 points)—enabling procedure-specific manipulation.

\subsection{Procedure-Specific Gaussian RBF Deformation}
\label{sec:deformation}

Given extracted landmarks $\mathbf{L}$, we apply procedure-specific
deformations using Gaussian radial basis function (RBF) interpolation.
Each procedure $p \in \{\text{rhinoplasty, blepharoplasty, rhytidectomy,
orthognathic}\}$ defines a set of \emph{control handles}
$\mathcal{H}_p = \{(k_j, \boldsymbol{\delta}_j, r_j)\}_{j=1}^{M_p}$,
where $k_j$ is the anchor landmark index, $\boldsymbol{\delta}_j \in
\RR^2$ is the displacement vector, and $r_j$ is the influence radius.

For a given intensity parameter $\alpha \in [0, 1]$, each landmark
$\mathbf{l}_i$ is displaced by the superposition of all active handles:

\begin{equation}
\label{eq:rbf}
\mathbf{l}_i' = \mathbf{l}_i + \sum_{j=1}^{M_p} \alpha \cdot
\boldsymbol{\delta}_j \cdot \exp\!\left(
-\frac{\|\mathbf{l}_i - \mathbf{l}_{k_j}\|^2}{2\,r_j^2}
\right)
\end{equation}

The Gaussian kernel ensures smooth, spatially localized deformation:
landmarks near a control handle receive large displacements while distant
landmarks remain unaffected, preventing anatomically implausible global
distortions.

Displacement vectors $\boldsymbol{\delta}_j$ are calibrated per procedure
from anthropometric studies~\cite{singh2010plastic}.  For example,
rhinoplasty applies three anatomically distinct sub-deformations:
(a)~bilateral alar narrowing—symmetric inward pull on nostril landmarks
(indices 240, 236, 460, 456 with $\delta_x = \pm 2.5\alpha$, $r = 0.6r_0$);
(b)~tip refinement—upward displacement on tip landmarks (indices 1, 2,
94 with $\delta_y = -2.0\alpha$, $r = 0.5r_0$); and (c)~dorsum
narrowing—bilateral squeeze on bridge landmarks ($\delta_x = \pm
1.5\alpha$, $r = 0.5r_0$).  All displacements are bilaterally symmetric
to maintain facial symmetry, consistent with standard surgical technique.

\subsection{Multi-Channel Conditioning Signal}
\label{sec:conditioning}

The primary conditioning signal is the \textbf{tessellation mesh}
$\mathbf{C} \in \RR^{512 \times 512 \times 3}$, rendered from the
deformed landmarks $\mathbf{L}' = \{\mathbf{l}_i'\}$.  This is the full
2,556-edge MediaPipe Face Mesh triangulation drawn on a black canvas:
gray wireframe edges (192, 192, 192) for the tessellation surface and
white contour edges (255, 255, 255) for key anatomical features (eyes,
eyebrows, lips, jawline).  Critically, this rendering matches the
training distribution of
CrucibleAI/ControlNetMediaPipeFace~\cite{crucibleai_controlnet},
allowing us to use pre-trained weights directly.

Two additional signals are computed but used in the compositing stage
rather than as ControlNet input:

\begin{itemize}
    \item \textbf{Auto-Canny edges}: edge detection applied to a static
    anatomical adjacency wireframe (not dynamic Delaunay, which produces
    triangle inversions under large displacements).  Adaptive thresholds
    ($0.66\tilde{v}$, $1.33\tilde{v}$) based on median pixel intensity
    $\tilde{v}$ ensure robustness across Fitzpatrick skin
    types~\cite{chardon1991ita}.

    \item \textbf{Feathered surgical mask}: a procedure-specific region
    derived from the convex hull of relevant landmarks, dilated by $d_p$
    pixels and Gaussian-feathered with $\sigma_p$.  Boundary noise
    (2--4\,px perturbation) prevents visible seams.  Parameters are
    per-procedure: rhinoplasty ($d=30$, $\sigma=15$), blepharoplasty
    ($d=15$, $\sigma=10$), rhytidectomy ($d=40$, $\sigma=20$),
    orthognathic ($d=35$, $\sigma=18$).  The mask governs the compositing
    stage (Sec.~\ref{sec:compositing}).
\end{itemize}

\subsection{ControlNet-Guided Latent Diffusion}
\label{sec:diffusion}

We employ Stable Diffusion 1.5~\cite{rombach2022ldm} as the base
generative model, augmented with ControlNet~\cite{zhang2023controlnet}
for spatial conditioning.  Specifically, we use the CrucibleAI
ControlNet variant pre-trained on MediaPipe face mesh
renderings~\cite{crucibleai_controlnet}, which accepts the tessellation
mesh as its conditioning input.

The denoising objective follows standard LDM formulation.  Given a latent
encoding $\mathbf{z}_0 = \mathcal{E}(I)$ of the input image via the VAE
encoder $\mathcal{E}$, noise $\boldsymbol{\epsilon} \sim
\mathcal{N}(\mathbf{0}, \mathbf{I})$ is added to produce
$\mathbf{z}_t$ at timestep $t$.  The U-Net $\epsilon_\theta$ predicts
the noise conditioned on the text prompt $c_\text{text}$ and the
ControlNet conditioning $\mathbf{C}$:

\begin{equation}
\label{eq:ldm}
\mathcal{L}_\text{diff} = \mathbb{E}_{\mathbf{z}_0, \boldsymbol{\epsilon},
t}\!\left[\|\boldsymbol{\epsilon} - \epsilon_\theta(\mathbf{z}_t, t,
c_\text{text}, \mathbf{C})\|^2\right]
\end{equation}

The ControlNet branch $\mathcal{F}$ creates a trainable copy of the
encoder blocks of $\epsilon_\theta$, connected via zero-initialized
convolution layers.  At each resolution level $l$, the ControlNet output
is added to the corresponding skip connection of the frozen U-Net:
$\mathbf{h}_l^{\text{out}} = \mathbf{h}_l^{\text{frozen}} + \lambda_c
\cdot \mathcal{F}_l(\mathbf{C})$, where $\lambda_c$ is the conditioning
scale.

Procedure-specific text prompts guide the diffusion process toward
clinically appropriate appearance (\eg, ``high quality professional
photo of a person's face, refined nose shape, smooth nasal bridge,
natural skin texture''), while a shared negative prompt suppresses common
diffusion artifacts (blurriness, distortion, extra features).

\subsection{Training Data Generation}
\label{sec:training_data}

Paired before/after surgical data is extremely scarce and ethically
constrained.  We generate synthetic training pairs from large-scale face
datasets (CelebA-HQ~\cite{karras2018celeba}, FFHQ~\cite{karras2019ffhq})
using the following procedure:

\begin{enumerate}
    \item Extract landmarks from a clean face image.
    \item Apply random procedure presets at random intensity
    $\alpha \sim \mathcal{U}(0.3, 0.9)$.
    \item Render the tessellation mesh conditioning from deformed landmarks.
    \item Generate a TPS-warped target image using thin-plate spline
    interpolation from original to deformed landmark positions, with
    subsampled control points ($\sim$80 from 478) for computational
    efficiency.
    \item Apply clinical degradation augmentation to the \emph{input only}
    (never the target), simulating the domain gap between curated and
    clinical photographs.
\end{enumerate}

The training target is the TPS-warped image, which approximates the
post-surgical geometry.  While TPS warps can exhibit visible
interpolation artifacts (``rubber face''), the ControlNet learns to
generate photorealistic faces that match the deformed landmark
configuration, using the TPS output as geometric guidance rather than
copying its texture artifacts.  The conditioning signal encodes the
deformed geometry via the tessellation mesh.

\paragraph{Clinical Augmentation Pipeline.}
To bridge the domain gap between curated datasets and clinical photographs,
we apply 3–5 augmentations per sample from a pool of 8 clinical
degradations: point-source lighting ($p=0.40$), color temperature jitter
($p=0.60$), green fluorescent cast ($p=0.25$), JPEG compression at
quality 40–85 ($p=0.30$), Gaussian sensor noise with $\sigma \in [5,25]$
($p=0.40$), barrel/pincushion lens distortion ($p=0.30$), motion blur
($p=0.20$), and vignetting ($p=0.25$).  Augmentations are applied to
input images only, preserving clean targets for the diffusion loss.

\subsection{Loss Function}
\label{sec:loss}

Training proceeds in two phases with a composite loss:
\begin{equation}
\label{eq:loss}
\mathcal{L}_\text{total} = \mathcal{L}_\text{diff}
+ \lambda_\ell \mathcal{L}_\text{landmark}
+ \lambda_\text{id} \mathcal{L}_\text{identity}
+ \lambda_p \mathcal{L}_\text{perceptual}
\end{equation}

\textbf{Phase A} (current, synthetic data): only $\mathcal{L}_\text{diff}$
(Eq.~\ref{eq:ldm}) is active.  Auxiliary losses are disabled because
penalizing photorealistic output against TPS-warped targets would be
counterproductive.

\textbf{Phase B} (planned, clinical data): all four terms enabled once
paired clinical data with verified outcomes becomes available.  The
auxiliary loss infrastructure is implemented but not yet activated.

\begin{itemize}
    \item $\mathcal{L}_\text{landmark}$: normalized mean error (NME)
    between predicted and target landmark positions, computed only within
    the surgical mask region and normalized by inter-ocular distance (IOD):
    \begin{equation}
    \mathcal{L}_\text{landmark} = \frac{1}{|\mathcal{M}|} \sum_{i \in \mathcal{M}}
    \frac{\|\mathbf{l}_i^\text{pred} - \mathbf{l}_i^\text{target}\|}{d_\text{IOD}}
    \end{equation}

    \item $\mathcal{L}_\text{identity}$: ArcFace~\cite{deng2019arcface}
    cosine similarity loss using 512-dimensional identity embeddings from
    InsightFace (buffalo\_l model).  Procedure-dependent cropping is
    applied: full face for blepharoplasty, upper two-thirds for
    rhinoplasty, upper three-quarters for rhytidectomy.  Disabled for
    orthognathic surgery where jaw changes legitimately alter facial
    identity features.

    \item $\mathcal{L}_\text{perceptual}$: LPIPS~\cite{zhang2018lpips}
    perceptual distance computed \emph{outside} the surgical mask only,
    ensuring non-surgical regions remain unchanged.
\end{itemize}

Default weights: $\lambda_\ell = 0.1$, $\lambda_\text{id} = 0.05$,
$\lambda_p = 0.1$.

\subsection{Identity-Preserving Mask Compositing}
\label{sec:compositing}

The final output is produced by compositing the diffusion output into the
original photograph using the feathered surgical mask.  To prevent color
shifts between the generated and original regions, we perform histogram
matching in CIE LAB space: for each channel $c \in \{L, a, b\}$, the
generated region's statistics $(\mu_c^g, \sigma_c^g)$ are matched to the
original's $(\mu_c^o, \sigma_c^o)$ within the mask:
\begin{equation}
\label{eq:color}
\hat{v}_c = \frac{(v_c - \mu_c^g) \cdot \sigma_c^o}{\sigma_c^g} + \mu_c^o
\end{equation}
The composited result $\hat{I}$ blends via the feathered mask $\mathbf{m}$:
\begin{equation}
\hat{I} = \mathbf{m} \odot \hat{I}_\text{gen} + (1 - \mathbf{m}) \odot I_\text{orig}
\end{equation}

This ensures pixel-perfect preservation of all non-surgical regions
(hair, ears, background, clothing) while seamlessly integrating the
synthesized surgical outcome.

%======================================================================
\section{Experimental Setup}
\label{sec:experiments}

\subsection{Datasets}

\paragraph{Training Data.}
We construct a synthetic training set of 200K pairs (50K per procedure)
from CelebA-HQ~\cite{karras2018celeba} (30K images) and an additional
14K diverse face images.  Real scraped before/after surgery images from
public forums supplement the synthetic data as a held-out validation
signal.

\paragraph{Evaluation Data.}
We evaluate on: (1) a held-out synthetic test set (1K pairs per procedure);
(2) the IIITD Plastic Surgery Face Database~\cite{singh2010plastic}
(900 subjects, 6 procedures, gated access); and (3) the HDA Plastic
Surgery Face Database (638 subjects).

\subsection{Baselines}

We compare against five baselines:
\begin{enumerate}
    \item \textbf{TPS-only}: thin-plate spline geometric warp without
    diffusion (our ablation baseline).
    \item \textbf{SD1.5 Img2Img}: Stable Diffusion 1.5 image-to-image
    with TPS-warped input and mask compositing, without ControlNet
    conditioning.
    \item \textbf{Pix2Pix}~\cite{isola2017pix2pix}: conditional GAN
    trained on our synthetic pairs.
    \item \textbf{CycleGAN}~\cite{zhu2017cyclegan}: unpaired
    image-to-image translation using scraped before/after data.
    \item \textbf{Crisalix-style morphing}: affine warp baseline using
    the same landmark displacements without generative modeling.
\end{enumerate}

\subsection{Evaluation Metrics}

We report five quantitative metrics:

\begin{itemize}
    \item \textbf{FID}~\cite{heusel2017fid}: Fr\'echet Inception Distance
    between generated and real post-surgery distributions (lower is better).
    \item \textbf{LPIPS}~\cite{zhang2018lpips}: perceptual distance
    between generated and target images (lower is better).
    \item \textbf{NME}: normalized mean landmark error between predicted
    and target landmark configurations, normalized by inter-ocular
    distance (lower is better).
    \item \textbf{SSIM}~\cite{wang2004ssim}: structural similarity index
    (higher is better; relaxed target $>0.80$).
    \item \textbf{ArcFace cosine similarity}: identity preservation score
    using ArcFace~\cite{deng2019arcface} embeddings (higher is better).
\end{itemize}

All metrics are stratified by Fitzpatrick skin type (I--VI) using the
Individual Typology Angle (ITA)~\cite{chardon1991ita} to assess equitable
performance across skin tones:
$\text{ITA} = \arctan\!\left(\frac{L^* - 50}{b^*}\right) \cdot
\frac{180}{\pi}$, where $L^*$ and $b^*$ are CIE LAB lightness and
yellow-blue chrominance.

\subsection{Ablation Studies}

We conduct ablations to evaluate:
\begin{enumerate}
    \item \textbf{Conditioning channels}: tessellation mesh alone vs.\
    mesh + Canny vs.\ full 3-channel (mesh + Canny + mask).
    \item \textbf{Deformation model}: random displacement vs.\ uniform
    vs.\ procedure-specific RBF (Eq.~\ref{eq:rbf}).
    \item \textbf{Compositing strategy}: direct output vs.\ mask
    compositing vs.\ mask + LAB color matching.
    \item \textbf{Clinical augmentation}: with vs.\ without the
    8-augmentation degradation pipeline.
    \item \textbf{Fine-tuning}: pre-trained CrucibleAI ControlNet vs.\
    fine-tuned on our surgical pairs.
\end{enumerate}

%======================================================================
\section{Results}
\label{sec:results}

\begin{table}[t]
\centering
\caption{Quantitative comparison across four surgical procedures.
Best results in \textbf{bold}, second-best \underline{underlined}.
Results are reported as mean $\pm$ std over 3 random seeds.}
\label{tab:results}
\small
\begin{tabular}{@{}lccccc@{}}
\toprule
Method & FID $\downarrow$ & LPIPS $\downarrow$ & NME $\downarrow$ & SSIM $\uparrow$ & ArcFace $\uparrow$ \\
\midrule
\multicolumn{6}{c}{\textit{Rhinoplasty}} \\
\midrule
TPS-only          & --    & --    & --    & --    & --   \\
SD1.5 Img2Img     & --    & --    & --    & --    & --   \\
Pix2Pix           & --    & --    & --    & --    & --   \\
Morphing          & --    & --    & --    & --    & --   \\
\textbf{\method{}} & --   & --    & --    & --    & --   \\
\midrule
\multicolumn{6}{c}{\textit{Blepharoplasty}} \\
\midrule
TPS-only          & --    & --    & --    & --    & --   \\
SD1.5 Img2Img     & --    & --    & --    & --    & --   \\
Pix2Pix           & --    & --    & --    & --    & --   \\
Morphing          & --    & --    & --    & --    & --   \\
\textbf{\method{}} & --   & --    & --    & --    & --   \\
\midrule
\multicolumn{6}{c}{\textit{Rhytidectomy}} \\
\midrule
TPS-only          & --    & --    & --    & --    & --   \\
SD1.5 Img2Img     & --    & --    & --    & --    & --   \\
Pix2Pix           & --    & --    & --    & --    & --   \\
Morphing          & --    & --    & --    & --    & --   \\
\textbf{\method{}} & --   & --    & --    & --    & --   \\
\midrule
\multicolumn{6}{c}{\textit{Orthognathic}} \\
\midrule
TPS-only          & --    & --    & --    & --    & --   \\
SD1.5 Img2Img     & --    & --    & --    & --    & --   \\
Pix2Pix           & --    & --    & --    & --    & --   \\
Morphing          & --    & --    & --    & --    & --   \\
\textbf{\method{}} & --   & --    & --    & --    & --   \\
\bottomrule
\end{tabular}
\end{table}

\begin{table}[t]
\centering
\caption{Ablation study on conditioning channels and compositing
strategy (rhinoplasty procedure, averaged over test set).}
\label{tab:ablation}
\small
\begin{tabular}{@{}lcccc@{}}
\toprule
Configuration & FID $\downarrow$ & LPIPS $\downarrow$ & NME $\downarrow$ & ArcFace $\uparrow$ \\
\midrule
Mesh only                   & -- & -- & -- & -- \\
Mesh + Canny                & -- & -- & -- & -- \\
Mesh + Canny + Mask (full)  & -- & -- & -- & -- \\
\midrule
Direct output (no composite)   & -- & -- & -- & -- \\
Mask composite (no color match) & -- & -- & -- & -- \\
Mask + LAB color match (full)   & -- & -- & -- & -- \\
\midrule
No clinical augmentation    & -- & -- & -- & -- \\
With clinical augmentation  & -- & -- & -- & -- \\
\bottomrule
\end{tabular}
\end{table}

\begin{table}[t]
\centering
\caption{Performance stratified by Fitzpatrick skin type (ITA-based
classification). FID and NME reported for rhinoplasty.}
\label{tab:fitzpatrick}
\small
\begin{tabular}{@{}lcccccc@{}}
\toprule
 & Type I & Type II & Type III & Type IV & Type V & Type VI \\
\midrule
FID $\downarrow$        & -- & -- & -- & -- & -- & -- \\
NME $\downarrow$        & -- & -- & -- & -- & -- & -- \\
ArcFace $\uparrow$      & -- & -- & -- & -- & -- & -- \\
\bottomrule
\end{tabular}
\end{table}

\paragraph{Quantitative Results.}
Tables~\ref{tab:results} and~\ref{tab:ablation} will be populated once
training completes.  We expect \method{} to outperform all baselines on
FID and LPIPS (photorealism), match TPS-only on NME (geometric accuracy
is determined by the deformation model, not the generative backbone),
and significantly outperform TPS-only and morphing baselines on perceptual
quality.

\paragraph{Qualitative Results.}
Figure~\ref{fig:qualitative} shows representative predictions across
all four procedures.  \method{} synthesizes realistic skin texture in
the modified region while preserving identity and non-surgical features.

%======================================================================
\section{Discussion and Limitations}
\label{sec:discussion}

\paragraph{Clinical Applicability.}
Because \method{} needs only a smartphone photograph, it could
realistically be deployed in outpatient clinics that lack 3D imaging.
The procedure-specific deformation presets give surgeons a meaningful
intensity dial, though the current interface exposes a single scalar
per procedure rather than independent sub-procedure controls (\eg,
separate alar narrowing and tip rotation sliders for rhinoplasty).
A finer-grained interface is straightforward to implement and would
likely improve clinical utility.

\paragraph{Limitations.}
Several limitations deserve candid acknowledgment.
First, our deformation model is strictly 2D and cannot capture
depth-dependent effects visible in profile views---dorsal hump reduction
in rhinoplasty side views being the most obvious example.  Coupling the
pipeline with FLAME~\cite{li2017flame} or another 3D morphable model
would address this and is a natural next step.
Second, the system predicts the \emph{healed} outcome; it does not model
scarring, swelling, or bruising, which are the immediate post-operative
reality.
Third, validation on real clinical data remains constrained by the
scarcity of publicly available paired surgical datasets with IRB
approval---a bottleneck shared by the entire field.
Finally, while we stratify evaluation by Fitzpatrick type, the training
corpora (CelebA-HQ, FFHQ) are well-documented to under-represent Types
V and VI~\cite{merler2019diversity}, and our results should be
interpreted with that imbalance in mind.

%======================================================================
\section{Conclusion}
\label{sec:conclusion}

We have shown that latent diffusion, conditioned on anatomically
deformed facial landmark meshes, can produce photorealistic surgical
outcome predictions from a single 2D photograph---without CT, without
structured light, and without paired clinical training data.  The key
ingredients are procedure-specific RBF deformation fields that encode
real surgical biomechanics, a conditioning scheme compatible with
pre-trained ControlNet weights, and a compositing stage that preserves
identity outside the surgical region.  Stratifying evaluation by
Fitzpatrick type keeps us honest about where the model works and where
it falls short.  Looking ahead, integration with 3D morphable models
would unlock profile-view prediction, and coupling with FEM-based
tissue simulation could ground the deformations in actual biomechanics
rather than learned proxies from anthropometric data.

%======================================================================
% FIGURES
%======================================================================

\begin{figure}[t]
\centering
% Placeholder for architecture diagram
\fbox{\parbox{0.95\textwidth}{\centering\vspace{2cm}
\textbf{Architecture Diagram Placeholder}\\[0.5em]
\small Input Photo $\rightarrow$ MediaPipe 478 Landmarks $\rightarrow$
Gaussian RBF Deformation $\rightarrow$ 3-Channel Conditioning\\
(Tessellation Mesh + Canny Edges + Feathered Mask) $\rightarrow$
ControlNet/SD1.5 $\rightarrow$ LAB Compositing $\rightarrow$ Output
\vspace{2cm}}}
\caption{Overview of the \method{} pipeline.  A single input photograph
is processed through five stages: (1) dense 478-point landmark extraction,
(2) procedure-specific RBF deformation, (3) three-channel conditioning
generation, (4) ControlNet-guided latent diffusion, and (5)
identity-preserving mask compositing with LAB color matching.}
\label{fig:architecture}
\end{figure}

\begin{figure}[t]
\centering
% Placeholder for qualitative results
\fbox{\parbox{0.95\textwidth}{\centering\vspace{3cm}
\textbf{Qualitative Results Placeholder}\\[0.5em]
\small 4 rows (one per procedure) $\times$ 5 columns:\\
Input $|$ Landmarks $|$ TPS Warp $|$ \method{} Output $|$ Mask Overlay
\vspace{3cm}}}
\caption{Qualitative results across four procedures.  For each row:
(a)~input photograph, (b)~extracted + deformed landmarks, (c)~TPS
geometric warp (baseline), (d)~\method{} output, (e)~surgical mask
overlay.  \method{} generates photorealistic texture in the surgical
region while the TPS baseline produces visible warping artifacts.}
\label{fig:qualitative}
\end{figure}

%======================================================================
% REFERENCES
%======================================================================

\bibliographystyle{splncs04}
\bibliography{references}

\end{document}
